\section{Introduction}

An LLM data agent turns a request into a structured plan, executable query, and released answer. Policy enters this path more than once: a capability manifest tells the model what it may request; a grammar defines valid plans; a validator accepts or rejects a plan; a compiler turns it into SQL; database policy constrains execution; and a release check decides what may leave the trusted boundary. We call these six components the \emph{policy pipeline}. They use different representations and are often deployed on different schedules.

The resulting failure is not ordinary component failure. Every component may pass its own tests while the transition between two components drops tenant identity, changes a metric definition, uses a stale alias, loses lineage, or releases a result under an obsolete rule. We call a mismatch between a canonical obligation and a policy-pipeline surface \emph{divergence}; when an update creates it, we call it \emph{policy drift}.

\paragraph{Worked example.}
Suppose analyst Alice belongs to tenant $A$ and may request tenant-$A$ ticket aggregates, but not tenant-$B$ rows or raw customer email. The canonical policy is version 7. The validator also runs v7, while the compiler still runs v5 and uses \texttt{legacy\_tenant\_id}. Alice asks for ``escalations by region.'' The validator accepts the semantic plan
\[
\langle\texttt{escalated\_tickets},[\texttt{region}],
\texttt{last\_month},100\rangle .
\]
Compiling that accepted plan is a \emph{lowering}: a translation from the typed plan into SQL. The stale compiler binds the tenant obligation to the wrong column. On a two-row database containing one distinguishing tenant-$B$ escalation, the emitted SQL reads both tenants. Database row-level security (RLS) may block the extra row, but that containment does not make the compiler correct.

\policystrata{} records the six deployed versions as a \emph{version vector}, here
\[
\begin{aligned}
\vec v={}&(v_{\rm manifest},v_{\rm grammar},v_{\rm validator},\\
          &v_{\rm compiler},v_{\rm db},v_{\rm release})
          =(7,7,7,5,7,7).
\end{aligned}
\]
It evaluates the canonical plan and the lowered query, checks each transition in order, and reports the compiler as the first violated contract. The witness retains Alice, the plan, the version vector, the distinguishing rows, the SQL, the RLS outcome, and the release decision. This example separates three facts that a final-answer assertion would collapse: the compiler violated tenant preservation, RLS contained the violation, and no external disclosure occurred.

The right oracle is therefore not literal equality between layers. A manifest may intentionally expose less than a validator accepts, and RLS may intentionally be stricter than an application check. \policystrata{} assigns each surface a responsibility and checks only the obligations that surface accepts or emits.

\paragraph{Scope and evidence boundary.}
\policystrata{} is a regression tester and scanner, not an authorization boundary. It assumes a human-reviewed canonical policy and trusted adapters that map implementation artifacts into the checking model. The v1 benchmark covers single-request, read-only analytics. Its deterministic generator selects plans from the same public operator taxonomy used by the simulator, so its perfect kill count measures internal coverage. A maintainer-operated deployment study verifies source-to-deployment binding and selected live denial behavior, but publishes only aggregate results, does not read customer rows, and lacks an authenticated synthetic principal. No evaluation observes an independently operated deployment or estimates recall on unknown production incidents.

\paragraph{Claims and non-claims.}
\begin{center}
\small
\begin{tabular}{@{}L{0.43\columnwidth}L{0.47\columnwidth}@{}}
\toprule
Claim & Non-claim \\
\midrule
Responsibility contracts detect and localize faults represented by the implemented model & The contracts do not prove a deployment's canonical policy is correct \\
The artifact covers its 22 operators and distinguishes gaps missed by two spec-derived comparators & 1720/1720 is not recall on unknown faults \\
The scanner can target real artifacts and an exact deployed revision & Deployment binding and denial probes are not authenticated cross-tenant validation \\
Every reported witness has a modeled contract violation in the checked state space & This is not a mechanized soundness proof or global completeness result \\
\bottomrule
\end{tabular}
\end{center}

\paragraph{Contributions.}
This paper contributes:

\begin{itemize}
  \item a model of a data agent as six non-equivalent policy surfaces with explicit exposure, validation, lowering, containment, and release contracts;
  \item a deterministic generator and replay checker that produces first-transition witnesses, plus counterfactual repair checks for attribution and bounded replay reduction for witness size;
  \item \bench, with 22 named operators, explicit suite provenance, equivalent/invalid accounting, freeze manifests, two spec-derived comparators, and separate synthetic, spec-blind, public-fault, brownfield, and historical-revision evidence;
  \item a documented retargeting contract plus an aggregate, deployment-linked read-only study whose private operational identifiers are withheld.
\end{itemize}
